\documentclass[../../main.tex]{subfiles}% 注意这里的文件路径不能用 ./main.tex ,否则用latexmk编译子文件会报错
\graphicspath{{\subfix{./image/}}{\subfix{../../image/}}} % 指定图片引用路径,后续可以直接使用图片文件名
% 如果图片存放在上级目录,则使用 ../../image/ 作为备用路径

% 例如:
% \begin{figure}[H]
% \centering
% \includegraphics[width=0.3\textwidth]{图.png}
% \caption{}
% \label{figure:图}
% \end{figure}
% 注意:上述\label{}一定要放在\caption{}之后,否则引用图片序号会只会显示??.

\begin{document}
% SOURCE p=103; printed=88; object=ch03-lebesgue-radon-nikodym-title; region=lower
\section{Lebesgue–Radon–Nikodym 定理}\label{section:ch03-lebesgue-radon-nikodym}
% SOURCE p=103; printed=88; object=ch03-absolute-continuity; region=lower
\begin{definition}[绝对连续]\label{definition:ch03-absolute-continuity}
设 $\nu$ 为符号测度，$\mu$ 为 $(X,\mathscr{M})$ 上的正测度。若对每个满足 $\mu(E)=0$ 的 $E\in\mathscr{M}$，均有 $\nu(E)=0$，则称 $\nu$ 关于 $\mu$ \textbf{绝对连续}，记作 $\nu\ll\mu$。
\end{definition}
% SOURCE p=103; printed=88; object=ch03-absolute-continuity-variations; region=bottom
\begin{proposition}\label{proposition:ch03-absolute-continuity-variations}
$\nu\ll\mu$ 当且仅当 $|\nu|\ll\mu$，也当且仅当 $\nu^+\ll\mu$ 且 $\nu^-\ll\mu$。
\end{proposition}
\begin{proof}
这容易验证（\cref{exercise:ch03-absolute-continuity-variations}）。
\end{proof}
% SOURCE p=103,104; printed=88,89; object=ch03-singular-absolute-continuity-zero; region=bottom/top
在某种意义上，绝对连续性与相互奇异性彼此相对。更准确地说：
\begin{proposition}\label{proposition:ch03-singular-absolute-continuity-zero}
若 $\nu\perp\mu$ 且 $\nu\ll\mu$，则 $\nu=0$。
\end{proposition}
\begin{proof}
若 $E,F$ 为彼此不交的集合，满足 $E\cup F=X$ 且 $\mu(E)=|\nu|(F)=0$，则由 $\nu\ll\mu$ 得 $|\nu|(E)=0$，从而 $|\nu|=0$，故 $\nu=0$。
\end{proof}
% SOURCE p=104; printed=89; object=ch03-absolute-continuity-signed-extension; region=top
\begin{definition}[关于符号测度的绝对连续性]\label{definition:ch03-absolute-continuity-signed-extension}
可将绝对连续性的概念推广到 $\mu$ 为符号测度的情形：规定 $\nu\ll\mu$ 当且仅当 $\nu\ll|\mu|$。
\end{definition}
\begin{remark}
本书不需要这个更一般的定义。
\end{remark}
% SOURCE p=104; printed=89; object=ch03-absolute-continuity-intro; region=upper
“绝对连续性”一词来自实变量理论，见\cref{section:ch03-bounded-variation}。对有限符号测度，它等价于另一个显然属于连续性形式的条件。
% SOURCE p=104; printed=89; object=ch03-absolute-continuity-epsilon-delta; number=3.5; region=upper
\begin{theorem}\label{theorem:ch03-absolute-continuity-epsilon-delta}
设 $\nu$ 为有限符号测度，$\mu$ 为 $(X,\mathscr{M})$ 上的正测度。则 $\nu\ll\mu$，当且仅当对每个 $\varepsilon>0$，都存在 $\delta>0$，使得只要 $\mu(E)<\delta$，就有 $|\nu(E)|<\varepsilon$。
\end{theorem}
\begin{proof}
由于 $\nu\ll\mu$ 当且仅当 $|\nu|\ll\mu$，且 $|\nu(E)|\le|\nu|(E)$，只需假设 $\nu=|\nu|$ 为正测度。显然，$\varepsilon$-$\delta$ 条件蕴含 $\nu\ll\mu$。另一方面，若 $\varepsilon$-$\delta$ 条件不成立，则存在 $\varepsilon>0$，使对每个 $n\in\mathbb{N}$，都能找到 $E_n\in\mathscr{M}$，满足 $\mu(E_n)<2^{-n}$ 且 $\nu(E_n)\ge\varepsilon$。令 $F_k=\bigcup_k^\infty E_n$，$F=\bigcap_1^\infty F_k$。则 $\mu(F_k)<\sum_k^\infty2^{-n}=2^{1-k}$，故 $\mu(F)=0$；但对所有 $k$ 都有 $\nu(F_k)\ge\varepsilon$，由 $\nu$ 有限可得 $\nu(F)=\lim\nu(F_k)\ge\varepsilon$。因此 $\nu\ll\mu$ 不成立。
\end{proof}
% SOURCE p=104; printed=89; object=ch03-density-absolute-continuity; region=middle
\begin{proposition}\label{proposition:ch03-density-absolute-continuity}
若 $\mu$ 为测度，$f$ 为广义 $\mu$-可积函数，则由 $\nu(E)=\int_E f\,\mathrm{d}\mu$ 定义的符号测度 $\nu$ 关于 $\mu$ 绝对连续，并且 $\nu$ 有限当且仅当 $f\in L^1(\mu)$。
\end{proposition}
对于任意复值函数 $f\in L^1(\mu)$，可将前面的定理分别用于 $\operatorname{Re}f$ 和 $\operatorname{Im}f$，得到下面这个有用的结论。
% SOURCE p=104; printed=89; object=ch03-integral-absolute-continuity; number=3.6; region=middle
\begin{corollary}\label{corollary:ch03-integral-absolute-continuity}
若 $f\in L^1(\mu)$，则对每个 $\varepsilon>0$，都存在 $\delta>0$，使得只要 $\mu(E)<\delta$，就有 $|\int_E f\,\mathrm{d}\mu|<\varepsilon$。
\end{corollary}
% SOURCE p=104; printed=89; object=ch03-density-notation; region=middle-lower
\begin{definition}[测度密度的记号]\label{definition:ch03-density-notation}
用记号
\begin{align*}
\mathrm{d}\nu=f\,\mathrm{d}\mu
\end{align*}
表示关系 $\nu(E)=\int_E f\,\mathrm{d}\mu$。有时也略微不严格地称它为“符号测度 $f\,\mathrm{d}\mu$”。
\end{definition}
% SOURCE p=104; printed=89; object=ch03-decomposition-intro; region=lower
现在来讨论本节的主要定理，它完整描述了一个符号测度相对于给定正测度的结构。先证明一个技术性引理。
% SOURCE p=104; printed=89; object=ch03-positive-domination; number=3.7; region=lower
\begin{lemma}\label{lemma:ch03-positive-domination}
设 $\nu,\mu$ 为 $(X,\mathscr{M})$ 上的有限测度。则或者 $\nu\perp\mu$，或者存在 $\varepsilon>0$ 和 $E\in\mathscr{M}$，使 $\mu(E)>0$ 且在 $E$ 上 $\nu\ge\varepsilon\mu$（即 $E$ 为 $\nu-\varepsilon\mu$ 的正集）。
\end{lemma}
\begin{proof}
设 $X=P_n\cup N_n$ 为 $\nu-n^{-1}\mu$ 的 Hahn 分解，令 $P=\bigcup_1^\infty P_n$，$N=\bigcap_1^\infty N_n=P^c$。则对每个 $n$，$N$ 都是 $\nu-n^{-1}\mu$ 的负集，即对所有 $n$，$0\le\nu(N)\le n^{-1}\mu(N)$，故 $\nu(N)=0$。若 $\mu(P)=0$，则 $\nu\perp\mu$。若 $\mu(P)>0$，则对某个 $n$ 有 $\mu(P_n)>0$，而 $P_n$ 为 $\nu-n^{-1}\mu$ 的正集。
\end{proof}
% SOURCE p=105; printed=90; object=ch03-lebesgue-radon-nikodym; number=3.8; region=top
\begin{theorem}[Lebesgue–Radon–Nikodym 定理]\label{theorem:ch03-lebesgue-radon-nikodym}
设 $\nu$ 为 $\sigma$-有限符号测度，$\mu$ 为 $(X,\mathscr{M})$ 上的 $\sigma$-有限正测度。则存在唯一的 $\sigma$-有限符号测度 $\lambda,\rho$，使
\begin{align*}
\lambda\perp\mu,\qquad \rho\ll\mu,\qquad \nu=\lambda+\rho.
\end{align*}
此外，存在广义 $\mu$-可积函数 $f:X\to\mathbb{R}$，使 $\mathrm{d}\rho=f\,\mathrm{d}\mu$，并且任意两个这样的函数都 $\mu$-几乎处处相等。
\end{theorem}
\begin{proof}
\textbf{情形一：}$\nu,\mu$ 为有限正测度。令
\begin{align*}
\mathscr{F}=\left\{f:X\to[0,\infty]:\int_E f\,\mathrm{d}\mu\le\nu(E)\text{ 对所有 }E\in\mathscr{M}\text{ 成立}\right\}.
\end{align*}
由于 $0\in\mathscr{F}$，$\mathscr{F}$ 非空。此外，若 $f,g\in\mathscr{F}$，则 $h=\max(f,g)\in\mathscr{F}$，因为令 $A=\{x:f(x)>g(x)\}$，对任意 $E\in\mathscr{M}$，有
\begin{align*}
\int_E h\,\mathrm{d}\mu
&=\int_{E\cap A}f\,\mathrm{d}\mu+\int_{E\setminus A}g\,\mathrm{d}\mu\\
&\le\nu(E\cap A)+\nu(E\setminus A)=\nu(E).
\end{align*}
令 $a=\sup\{\int f\,\mathrm{d}\mu:f\in\mathscr{F}\}$，注意到 $a\le\nu(X)<\infty$，选取序列 $\{f_n\}\subseteq\mathscr{F}$，使 $\int f_n\,\mathrm{d}\mu\to a$。令 $g_n=\max(f_1,\ldots,f_n)$，$f=\sup_n f_n$。则 $g_n\in\mathscr{F}$，$g_n$ 逐点递增趋于 $f$，且 $\int g_n\,\mathrm{d}\mu\ge\int f_n\,\mathrm{d}\mu$。由此 $\lim\int g_n\,\mathrm{d}\mu=a$，再由单调收敛定理，$f\in\mathscr{F}$ 且 $\int f\,\mathrm{d}\mu=a$。（特别地，$f<\infty$ 几乎处处成立，所以可以取 $f$ 处处为实值。）

我们断言测度 $\mathrm{d}\lambda=\mathrm{d}\nu-f\,\mathrm{d}\mu$（由 $f\in\mathscr{F}$，它为正测度）关于 $\mu$ 奇异。否则，由\cref{lemma:ch03-positive-domination}，存在 $E\in\mathscr{M}$ 和 $\varepsilon>0$，使 $\mu(E)>0$ 且在 $E$ 上 $\lambda\ge\varepsilon\mu$。于是 $\varepsilon\chi_E\,\mathrm{d}\mu\le\mathrm{d}\lambda=\mathrm{d}\nu-f\,\mathrm{d}\mu$，即 $(f+\varepsilon\chi_E)\,\mathrm{d}\mu\le\mathrm{d}\nu$。所以 $f+\varepsilon\chi_E\in\mathscr{F}$，且 $\int(f+\varepsilon\chi_E)\,\mathrm{d}\mu=a+\varepsilon\mu(E)>a$，这与 $a$ 的定义矛盾。

这样就证明了 $\lambda,f$ 及 $\mathrm{d}\rho=f\,\mathrm{d}\mu$ 的存在性。至于唯一性，若还有 $\mathrm{d}\nu=\mathrm{d}\lambda'+f'\,\mathrm{d}\mu$，则 $\mathrm{d}\lambda-\mathrm{d}\lambda'=(f'-f)\,\mathrm{d}\mu$。但 $\lambda-\lambda'\perp\mu$（见\cref{exercise:ch03-countable-singular-absolute}），而 $(f'-f)\,\mathrm{d}\mu\ll\mathrm{d}\mu$，所以 $\mathrm{d}\lambda-\mathrm{d}\lambda'=(f'-f)\,\mathrm{d}\mu=0$，从而 $\lambda=\lambda'$，并且由\cref{proposition:ch02-lone-ae-equivalence}，$f=f'$，$\mu$-几乎处处成立。$\mu,\nu$ 为有限测度的情形得证。

\textbf{情形二：}$\mu,\nu$ 为 $\sigma$-有限测度。$X$ 可以写成可数个 $\mu$-有限集的不交并，也可以写成可数个 $\nu$-有限集的不交并。对这些集合取交，得到 $\mathscr{M}$ 中的不交序列 $\{A_j\}$，使对每个 $j$，$\mu(A_j),\nu(A_j)$ 都有限，且 $X=\bigcup_1^\infty A_j$。定义 $\mu_j(E)=\mu(E\cap A_j)$ 和 $\nu_j(E)=\nu(E\cap A_j)$。由上面的推理，对每个 $j$ 有 $\mathrm{d}\nu_j=\mathrm{d}\lambda_j+f_j\,\mathrm{d}\mu_j$，其中 $\lambda_j\perp\mu_j$。由于 $\mu_j(A_j^c)=\nu_j(A_j^c)=0$，有 $\lambda_j(A_j^c)=\nu_j(A_j^c)-\int_{A_j^c}f\,\mathrm{d}\mu_j=0$，并可假设 $f_j$ 在 $A_j^c$ 上为零。令 $\lambda=\sum_1^\infty\lambda_j$，$f=\sum_1^\infty f_j$。则 $\mathrm{d}\nu=\mathrm{d}\lambda+f\,\mathrm{d}\mu$，$\lambda\perp\mu$（见\cref{exercise:ch03-countable-singular-absolute}），且 $\mathrm{d}\lambda$ 与 $f\,\mathrm{d}\mu$ 均为 $\sigma$-有限，符合要求。唯一性同前。

\textbf{一般情形：}若 $\nu$ 为符号测度，将上述论证用于 $\nu^+$ 和 $\nu^-$，再将所得结果相减即可。
\end{proof}
% SOURCE p=106; printed=91; object=ch03-lebesgue-decomposition; region=top
\begin{definition}[Lebesgue 分解]\label{definition:ch03-lebesgue-decomposition}
满足 $\lambda\perp\mu$、$\rho\ll\mu$ 的分解 $\nu=\lambda+\rho$，称为 $\nu$ 关于 $\mu$ 的\textbf{Lebesgue 分解}。
\end{definition}
% SOURCE p=106; printed=91; object=ch03-radon-nikodym; region=top
\begin{corollary}[Radon–Nikodym 定理]\label{corollary:ch03-radon-nikodym}
在\cref{theorem:ch03-lebesgue-radon-nikodym}中，若 $\nu\ll\mu$，则对某个 $f$，有 $\mathrm{d}\nu=f\,\mathrm{d}\mu$。
\end{corollary}
% SOURCE p=106; printed=91; object=ch03-radon-nikodym-derivative; region=upper
\begin{definition}[Radon–Nikodym 导数]\label{definition:ch03-radon-nikodym-derivative}
上述函数 $f$ 称为 $\nu$ 关于 $\mu$ 的\textbf{Radon–Nikodym 导数}，记作 $\mathrm{d}\nu/\mathrm{d}\mu$，于是
\begin{align*}
\mathrm{d}\nu=\frac{\mathrm{d}\nu}{\mathrm{d}\mu}\,\mathrm{d}\mu.
\end{align*}
严格地说，$\mathrm{d}\nu/\mathrm{d}\mu$ 应理解为与 $f$ 在 $\mu$-几乎处处意义下相等的函数所组成的等价类。
\end{definition}
% SOURCE p=106; printed=91; object=ch03-radon-nikodym-linearity; region=upper
微分记号 $\mathrm{d}\nu/\mathrm{d}\mu$ 所暗示的公式一般是正确的。例如，显然有
\begin{proposition}\label{proposition:ch03-radon-nikodym-linearity}
\begin{align*}
\frac{\mathrm{d}(\nu_1+\nu_2)}{\mathrm{d}\mu}=\frac{\mathrm{d}\nu_1}{\mathrm{d}\mu}+\frac{\mathrm{d}\nu_2}{\mathrm{d}\mu}.
\end{align*}
\end{proposition}
并且有如下链式法则。
% SOURCE p=106; printed=91; object=ch03-radon-nikodym-chain; number=3.9; region=middle
\begin{proposition}\label{proposition:ch03-radon-nikodym-chain}
设 $\nu$ 为 $\sigma$-有限符号测度，$\mu,\lambda$ 为 $(X,\mathscr{M})$ 上的 $\sigma$-有限测度，且 $\nu\ll\mu$、$\mu\ll\lambda$。
\begin{enumerate}
\item\label{item:ch03-radon-nikodym-integration} 若 $g\in L^1(\nu)$，则 $g(\mathrm{d}\nu/\mathrm{d}\mu)\in L^1(\mu)$，且
\begin{align*}
\int g\,\mathrm{d}\nu=\int g\frac{\mathrm{d}\nu}{\mathrm{d}\mu}\,\mathrm{d}\mu.
\end{align*}
\item\label{item:ch03-radon-nikodym-chain} 有 $\nu\ll\lambda$，且
\begin{align*}
\frac{\mathrm{d}\nu}{\mathrm{d}\lambda}=\frac{\mathrm{d}\nu}{\mathrm{d}\mu}\frac{\mathrm{d}\mu}{\mathrm{d}\lambda}\qquad\lambda\text{-几乎处处}.
\end{align*}
\end{enumerate}
\end{proposition}
\begin{proof}
分别考虑 $\nu^+$ 和 $\nu^-$，可以假设 $\nu\ge0$。由 $\mathrm{d}\nu/\mathrm{d}\mu$ 的定义，等式 $\int g\,\mathrm{d}\nu=\int g(\mathrm{d}\nu/\mathrm{d}\mu)\,\mathrm{d}\mu$ 在 $g=\chi_E$ 时成立。由线性，它对简单函数成立；再由单调收敛定理，它对非负可测函数成立；最后再次利用线性，得到它对 $L^1(\nu)$ 中的函数成立。分别以 $\mu,\lambda$ 代替 $\nu,\mu$，并令 $g=\chi_E(\mathrm{d}\nu/\mathrm{d}\mu)$，得到
\begin{align*}
\nu(E)=\int_E\frac{\mathrm{d}\nu}{\mathrm{d}\mu}\,\mathrm{d}\mu
=\int_E\frac{\mathrm{d}\nu}{\mathrm{d}\mu}\frac{\mathrm{d}\mu}{\mathrm{d}\lambda}\,\mathrm{d}\lambda
\end{align*}
对所有 $E\in\mathscr{M}$ 成立。由\cref{proposition:ch02-lone-ae-equivalence}，$(\mathrm{d}\nu/\mathrm{d}\lambda)=(\mathrm{d}\nu/\mathrm{d}\mu)(\mathrm{d}\mu/\mathrm{d}\lambda)$，$\lambda$-几乎处处成立。
\end{proof}
% SOURCE p=106; printed=91; object=ch03-radon-nikodym-reciprocal; number=3.10; region=lower
\begin{corollary}\label{corollary:ch03-radon-nikodym-reciprocal}
若 $\mu\ll\lambda$ 且 $\lambda\ll\mu$，则 $(\mathrm{d}\lambda/\mathrm{d}\mu)(\mathrm{d}\mu/\mathrm{d}\lambda)=1$，几乎处处成立（关于 $\lambda$ 或 $\mu$ 都如此）。
\end{corollary}
% SOURCE p=106; printed=91; object=ch03-dirac-nonexample; region=lower
\begin{example}[不存在 Radon–Nikodym 导数的例子]\label{example:ch03-dirac-nonexample}
设 $\mu$ 为 $(\mathbb{R},\mathscr{B}_{\mathbb{R}})$ 上的 Lebesgue 测度，$\nu$ 为原点处的点质量。显然 $\nu\perp\mu$。这个并不存在的 Radon–Nikodym 导数 $\mathrm{d}\nu/\mathrm{d}\mu$，通常被称为 Dirac $\delta$-函数。
\end{example}
% SOURCE p=106; printed=91; object=ch03-common-dominating-measure-intro; region=lower
以一个简单但重要的观察结束本节。
% SOURCE p=106; printed=91; object=ch03-common-dominating-measure; number=3.11; region=bottom
\begin{proposition}\label{proposition:ch03-common-dominating-measure}
若 $\mu_1,\ldots,\mu_n$ 为 $(X,\mathscr{M})$ 上的测度，则存在测度 $\mu$，使对所有 $j$ 都有 $\mu_j\ll\mu$；取 $\mu=\sum_1^n\mu_j$ 即可。
\end{proposition}
\begin{proof}
证明是显然的。
\end{proof}
% SOURCE p=107; printed=92; object=ch03-lrn-exercises-title; region=top
\subsection*{习题}
% SOURCE p=107; printed=92; object=ch03-ex08; number=8; region=top
\begin{exercise}\label{exercise:ch03-absolute-continuity-variations}
$\nu\ll\mu$ 当且仅当 $|\nu|\ll\mu$，也当且仅当 $\nu^+\ll\mu$ 且 $\nu^-\ll\mu$。
\end{exercise}
% SOURCE p=107; printed=92; object=ch03-ex09; number=9; region=top
\begin{exercise}\label{exercise:ch03-countable-singular-absolute}
设 $\{\nu_j\}$ 为正测度序列。若对所有 $j$，$\nu_j\perp\mu$，则 $\sum_1^\infty\nu_j\perp\mu$；若对所有 $j$，$\nu_j\ll\mu$，则 $\sum_1^\infty\nu_j\ll\mu$。
\end{exercise}
% SOURCE p=107; printed=92; object=ch03-ex10; number=10; region=upper
\begin{exercise}\label{exercise:ch03-nonfinite-absolute-continuity}
当 $\nu$ 不有限时，\cref{theorem:ch03-absolute-continuity-epsilon-delta}可能不成立。（考虑 $(0,1)$ 上的 $\mathrm{d}\nu(x)=\mathrm{d}x/x$ 和 $\mathrm{d}\mu(x)=\mathrm{d}x$；或者 $\mathbb{N}$ 上的 $\nu=$ 计数测度，$\mu(E)=\sum_{n\in E}2^{-n}$。）
\end{exercise}
% SOURCE p=107; printed=92; object=ch03-ex11; number=11; region=upper
\begin{exercise}\label{exercise:ch03-uniform-integrability}
设 $\mu$ 为正测度。
\begin{definition*}[一致可积]
函数族 $\{f_\alpha\}_{\alpha\in A}\subseteq L^1(\mu)$ 称为\textbf{一致可积}，若对每个 $\varepsilon>0$，都存在 $\delta>0$，使得只要 $\mu(E)<\delta$，就对所有 $\alpha\in A$ 有 $|\int_E f_\alpha\,\mathrm{d}\mu|<\varepsilon$。
\end{definition*}
\begin{enumerate}
\item $L^1(\mu)$ 的每个有限子集都是一致可积的。
\item 若 $L^1(\mu)$ 中的序列 $\{f_n\}$ 在 $L^1$ 度量下收敛于 $f\in L^1(\mu)$，则 $\{f_n\}$ 一致可积。
\end{enumerate}
\end{exercise}
% SOURCE p=107; printed=92; object=ch03-ex12; number=12; region=upper-middle
\begin{exercise}\label{exercise:ch03-radon-nikodym-product}
对 $j=1,2$，设 $\mu_j,\nu_j$ 为 $(X_j,\mathscr{M}_j)$ 上的 $\sigma$-有限测度，且 $\nu_j\ll\mu_j$。则 $\nu_1\times\nu_2\ll\mu_1\times\mu_2$，且
\begin{align*}
\frac{\mathrm{d}(\nu_1\times\nu_2)}{\mathrm{d}(\mu_1\times\mu_2)}(x_1,x_2)
=\frac{\mathrm{d}\nu_1}{\mathrm{d}\mu_1}(x_1)\frac{\mathrm{d}\nu_2}{\mathrm{d}\mu_2}(x_2).
\end{align*}
\end{exercise}
% SOURCE p=107; printed=92; object=ch03-ex13; number=13; region=middle
\begin{exercise}\label{exercise:ch03-non-sigma-finite-lrn}
设 $X=[0,1]$，$\mathscr{M}=\mathscr{B}_{[0,1]}$，$m=$ Lebesgue 测度，$\mu=$ $\mathscr{M}$ 上的计数测度。
\begin{enumerate}
\item $m\ll\mu$，但对任意 $f$，均有 $\mathrm{d}m\ne f\,\mathrm{d}\mu$。
\item $\mu$ 关于 $m$ 没有 Lebesgue 分解。
\end{enumerate}
\end{exercise}
% SOURCE p=107; printed=92; object=ch03-ex14; number=14; region=middle
\begin{exercise}\label{exercise:ch03-extended-radon-nikodym}
若 $\nu$ 为任意符号测度，$\mu$ 为 $(X,\mathscr{M})$ 上的 $\sigma$-有限测度，且 $\nu\ll\mu$，则存在广义 $\mu$-可积函数 $f:X\to[-\infty,\infty]$，使 $\mathrm{d}\nu=f\,\mathrm{d}\mu$。提示：
\begin{enumerate}
\item 只需假设 $\mu$ 有限且 $\nu$ 为正测度。
\item 在这些假设下，存在关于 $\nu$ 为 $\sigma$-有限的集合 $E\in\mathscr{M}$，使对所有关于 $\nu$ 为 $\sigma$-有限的集合 $F$，都有 $\mu(E)\ge\mu(F)$。
\item Radon–Nikodym 定理在 $E$ 上适用。若 $F\cap E=\varnothing$，则或者 $\nu(F)=\mu(F)=0$，或者 $\mu(F)>0$ 且 $|\nu(F)|=\infty$。
\end{enumerate}
\end{exercise}
% SOURCE p=107; printed=92; object=ch03-ex15; number=15; region=lower
\begin{exercise}\label{exercise:ch03-decomposable-measure}
\begin{definition*}[可分解测度]
$(X,\mathscr{M})$ 上的测度 $\mu$ 称为\textbf{可分解的}，若存在集合族 $\mathscr{F}\subseteq\mathscr{M}$，具有以下性质：
\begin{enumerate}[label=(\roman*)]
\item 对所有 $F\in\mathscr{F}$，$\mu(F)<\infty$。
\item $\mathscr{F}$ 的成员两两不交，且其并为 $X$。
\item 若 $\mu(E)<\infty$，则 $\mu(E)=\sum_{F\in\mathscr{F}}\mu(E\cap F)$。
\item 若 $E\subseteq X$，且对所有 $F\in\mathscr{F}$，$E\cap F\in\mathscr{M}$，则 $E\in\mathscr{M}$。
\end{enumerate}
\end{definition*}
\begin{enumerate}
\item 每个 $\sigma$-有限测度都是可分解的。
\item 若 $\mu$ 可分解，$\nu$ 为 $(X,\mathscr{M})$ 上任意满足 $\nu\ll\mu$ 的符号测度，则存在可测函数 $f:X\to[-\infty,\infty]$，使对任意关于 $\mu$ 为 $\sigma$-有限的 $E$，都有 $\nu(E)=\int_E f\,\mathrm{d}\mu$，且在任意关于 $\nu$ 为 $\sigma$-有限的 $F\in\mathscr{F}$ 上，$|f|<\infty$。（若 $\nu$ 不是 $\sigma$-有限，使用\cref{exercise:ch03-extended-radon-nikodym}。）
\end{enumerate}
\end{exercise}
% SOURCE p=107; printed=92; object=ch03-ex16; number=16; region=bottom
\begin{exercise}\label{exercise:ch03-radon-nikodym-sum}
设 $\mu,\nu$ 为 $(X,\mathscr{M})$ 上的 $\sigma$-有限测度，且 $\nu\ll\mu$，令 $\lambda=\mu+\nu$。若 $f=\mathrm{d}\nu/\mathrm{d}\lambda$，则 $0\le f<1$，$\mu$-几乎处处成立，且 $\mathrm{d}\nu/\mathrm{d}\mu=f/(1-f)$。
\end{exercise}
% SOURCE p=108; printed=93; object=ch03-ex17; number=17; region=top
\begin{exercise}\label{exercise:ch03-conditional-expectation}
设 $(X,\mathscr{M},\mu)$ 为有限测度空间，$\mathscr{N}$ 为 $\mathscr{M}$ 的子 $\sigma$-代数，$\nu=\mu|_{\mathscr{N}}$。若 $f\in L^1(\mu)$，则存在 $g\in L^1(\nu)$（因此 $g$ 为 $\mathscr{N}$-可测），使对所有 $E\in\mathscr{N}$ 都有 $\int_E f\,\mathrm{d}\mu=\int_E g\,\mathrm{d}\nu$；若 $g'$ 是另一个这样的函数，则 $g=g'$，$\nu$-几乎处处成立。
\begin{definition*}[条件期望]
在概率论中，$g$ 称为 $f$ 在 $\mathscr{N}$ 上的\textbf{条件期望}。
\end{definition*}
\end{exercise}
\end{document}
